\section{Modelagem}
\label{sec:modelagem}

\subsection{Como ler o diagrama}

O modelo conceitual segue a notação de Peter Chen \cite{chen1976}, exigida pelo enunciado, na forma apresentada por Heuser \cite{heuser2009}, e foi desenhado no brModelo. A Figura~\ref{fig:legenda} reúne os elementos que aparecem no diagrama, com as mesmas cores e formas da exportação, e a Tabela~\ref{tab:cores} explica como ler cada um.

\begin{figure}[h!]
\centering
\begin{tikzpicture}[x=1cm,y=1cm]
\node[ent,fill=brforte] (f) at (0,0) {Pessoa};
\node[font=\sffamily\scriptsize,below=1pt] at (f.south) {entidade forte};
\node[ent,fill=brfraca] (w) at (3.2,0) {Edição};
\node[font=\sffamily\scriptsize,below=1pt] at (w.south) {entidade fraca};
\node[ent,fill=brassoc,minimum width=2.3cm,minimum height=0.9cm] (a) at (6.6,0) {};
\node[rel,fill=white,minimum width=1.6cm] at (a) {Preço};
\node[font=\sffamily\scriptsize,below=1pt] at (a.south) {entidade associativa};
\node[rel,fill=brauto] (g) at (10.2,0) {Agrega};
\node[font=\sffamily\scriptsize,below=1pt] at (g.south) {autorrelacionamento};
\node[rel,fill=brrel] (r) at (0,-1.9) {Realiza};
\node[font=\sffamily\scriptsize,below=1pt] at (r.south) {relacionamento};
\node[idt] (i) at (2.5,-1.75) {};
\node[font=\footnotesize,right=2pt] at (i.east) {\texttt{cp\_id\_pessoa}};
\node[font=\sffamily\scriptsize] at (3.25,-2.35) {atributo identificador};
\node[atr] (c) at (5.7,-1.75) {};
\node[font=\footnotesize,right=2pt] at (c.east) {\texttt{nm\_pessoa}};
\node[font=\sffamily\scriptsize] at (6.2,-2.35) {atributo};
\node[font=\footnotesize] (k) at (9.6,-1.75) {(mín, máx)};
\node[font=\sffamily\scriptsize] at (9.6,-2.35) {cardinalidade};
\end{tikzpicture}
\caption{Elementos do diagrama, com as cores e formas usadas no brModelo}
\label{fig:legenda}
\end{figure}

{\small
\begin{longtable}{@{}L{3.2cm} L{2.5cm} L{4.3cm} L{3.5cm}@{}}
\caption{Dicionário de cores e formas do diagrama}\label{tab:cores}\\
\toprule
Elemento & Como aparece & Como ler & Itens no modelo\\
\midrule
\endfirsthead
\toprule
Elemento & Como aparece & Como ler & Itens no modelo\\
\midrule
\endhead
\bottomrule
\endfoot
Entidade forte & retângulo azul & tem identidade própria e existe sem depender de outra entidade & Pessoa, Evento, Instituição, Inscrição, Termo, Certificado, Auditoria\\
Entidade fraca & retângulo amarelo & só existe ligada à entidade dona, por um relacionamento identificador com cardinalidade (1,1) do lado da dona & Edição, Categoria de Inscrição, Faixa de Inscrição, Atividade, Trabalho Científico, Período de Associação, Benefício, Documento de Comprovação, Pagamento, Item Financeiro\\
Entidade associativa & retângulo rosa envolvendo um losango & relacionamento N:N que ganhou atributos próprios e passou a ser tratado como entidade, podendo participar de outros relacionamentos & Preço, Vínculo Institucional, Inscrição em Atividade\\
Autorrelacionamento & losango verde & relacionamento de uma entidade com ela mesma. As duas linhas saem da mesma entidade, cada uma com um papel & Agrega (edição agregadora e edições agregadas)\\
Relacionamento & losango branco & liga as entidades das pontas. Atributos ligados ao losango pertencem à ligação, e não a uma das entidades & os demais 26 relacionamentos\\
Atributo identificador & círculo cheio & identifica cada instância da entidade & um por entidade (\cod{cp_id_*})\\
Atributo & círculo vazio & característica da entidade ou do relacionamento a que está ligado & todos os demais\\
Cardinalidade & par (mín, máx) junto à entidade & lida a partir do outro lado, como explicado abaixo & todas as ligações\\
\end{longtable}
}

O par escrito junto a uma entidade diz com quantas instâncias dela uma instância do outro lado se relaciona. O mínimo 0 indica participação opcional e o mínimo 1, obrigatória. O máximo 1 ou n indica se a ligação é única ou múltipla. Em Evento (1,1) Realiza (1,n) Edição, cada edição pertence a exatamente um evento e cada evento tem uma ou mais edições. As cores servem para facilitar a leitura. O tipo de cada entidade também está marcado nas propriedades do brModelo (entidade fraca, entidade associativa e autorrelacionamento).

\subsection{Evolução do modelo}

O modelo passou por três versões antes da final.

A primeira versão (Figura~\ref{fig:merv1}) tinha 34 entidades. Ela foi útil para levantar tudo o que o caso pedia, mas misturava parâmetros e subtipos com entidades e representava as ligações como atributos de chave estrangeira, sem losangos.

\begin{figure}[h!]
\centering
\mervum
\caption{Primeira versão do modelo conceitual, com 34 entidades e ligações desenhadas como atributos}
\label{fig:merv1}
\end{figure}

Depois do retorno do professor, a segunda versão existiu como dicionário: o grupo revisou cada entidade e cada ligação por escrito antes de redesenhar (Seção~\ref{sec:escopo}). A terceira versão levou esse dicionário para o brModelo, já com 20 entidades e os relacionamentos com cardinalidade. Na revisão dela, o grupo fez os ajustes que resultaram na versão final: o relacionamento entre Pessoa e Certificado ganhou o nome Certifica, porque o nome Recebe já era usado entre Instituição e Benefício. Foram criados Tarifa, que liga a atividade selecionada ao item cobrado, e Atesta, que liga o certificado de atividade à atividade. A cardinalidade de Publica do lado da Edição passou a ser (1,1). E todos os atributos foram desenhados no diagrama, conferidos um a um com as colunas do banco.

\subsection{Modelo conceitual final}

A Figura~\ref{fig:merfinal} mostra o diagrama completo: 20 entidades (7 fortes, 10 fracas e 3 associativas), 30 relacionamentos e todos os atributos. Como o diagrama é grande, as Figuras~\ref{fig:merA} a~\ref{fig:merF} repetem cada região ampliada, com o comentário das decisões de cada parte. As figuras são vetoriais, recortadas da exportação do brModelo, e podem ser ampliadas no PDF sem perda de nitidez. Como as regiões do diagrama se tocam, um recorte pode mostrar parte de outro. Cada elemento é comentado na parte em que aparece por inteiro.

\begin{landscape}
\begin{figure}[p]
\centering
\mergeral
\caption{Modelo conceitual final do EventOS}
\label{fig:merfinal}
\end{figure}
\end{landscape}

\subsubsection*{Parte A: evento e edição}

Evento é forte e Edição é fraca, identificada pelo evento (Realiza). Cada edição guarda o período, o local, a modalidade e as políticas que valem para ela: prazo de reserva do preço, regra de reembolso, exigência de credenciamento, frequência mínima para certificado e limite de trabalhos por inscrição (RF1, RF15, RF20, RF22, RE2, RE9). O autorrelacionamento Agrega liga uma edição agregadora às edições que acontecem junto com ela (RF2). Limitamos a agregação a um nível, para que a regra de acesso fique simples de verificar. Organiza liga Pessoa e Edição e guarda, em \cod{tp_papel}, o papel de cada membro da equipe (RF24).

\begin{figure}[h!]
\centering
\merA
\caption{Parte A do modelo conceitual: evento, edição e organização}
\label{fig:merA}
\end{figure}

\subsubsection*{Parte B: oferta da edição}

Categoria de Inscrição, Faixa de Inscrição e Atividade são fracas de Edição (Oferece, Estabelece e Programa), porque o enunciado diz que cada edição define as suas e que nomes iguais em edições diferentes não são a mesma coisa (RF5). Preço é a entidade associativa entre Categoria e Faixa. Cada preço combina uma categoria, uma faixa e um tipo de associação, e guarda o valor da inscrição e, no combo, o valor da associação (RF6, RF7). O relacionamento Restringe diz quais categorias podem selecionar uma atividade. Atividade sem restrição fica aberta a todas (RE17).

\begin{figure}[h!]
\centering
\merB
\caption{Parte B do modelo conceitual: categorias, faixas, preços e atividades}
\label{fig:merB}
\end{figure}

\subsubsection*{Parte C: pessoas e instituições}

Pessoa é o cadastro único de participantes, autores, organizadores e analistas. Período de Associação é fraca de Pessoa (Associa-se): cada período tem início e fim, e a situação em uma data é derivada deles (RF4). Vínculo Institucional é associativa entre Pessoa e Instituição, com tipo e período, porque a mesma pessoa pode voltar à mesma instituição em outro momento (RF16, RE14).

Benefício é fraco de Instituição (Recebe). O uso do benefício não virou entidade: o relacionamento Facilita liga o benefício à inscrição que o usou, e a data de uso fica em \cod{dt_uso_beneficio}, na inscrição, porque cada inscrição usa no máximo um benefício (RF17). Os losangos Efetua, Analisa e Assina, que também aparecem nesta região, ligam Pessoa à inscrição, ao documento analisado e ao trabalho, e são comentados nas Partes D e E.

\begin{figure}[h!]
\centering
\merC
\caption{Parte C do modelo conceitual: pessoas, associação, instituições e benefícios}
\label{fig:merC}
\end{figure}

\subsubsection*{Parte D: inscrição e financeiro}

Inscrição é forte, com identidade própria, porque a mesma pessoa pode ter mais de uma inscrição na mesma edição (por exemplo, uma expirada e uma nova). Ela se liga à pessoa (Efetua) e ao preço (Enquadra) e, pelo preço, à categoria, à faixa e à edição. Por isso a inscrição não guarda categoria nem edição.

Item Financeiro, Pagamento e Documento de Comprovação são fracos de Inscrição (Compõe, Transaciona e Garante). Cada item é uma parte do valor devido e pode vir de uma atividade (Tarifa) ou de um trabalho excedente (Sobretaxa) (RF12). Pagamento guarda cada tentativa, com forma, valor e situação (RF13, RF14). O reembolso fica no próprio pagamento (RE15). Documento guarda cada envio, e Analisa liga o documento à pessoa que o analisou (RF8, RE12).

Inscrição em Atividade é associativa entre Inscrição e Atividade, com data de inclusão, situação, prazo da reserva, cancelamento e presença (RF10, RF21, RE4). Origina liga a inscrição de combo ao período de associação criado quando ela é paga (RF7).

\begin{figure}[h!]
\centering
\merD
\caption{Parte D do modelo conceitual: inscrição, itens, pagamentos, documentos e atividades selecionadas}
\label{fig:merD}
\end{figure}

\subsubsection*{Parte E: publicação e termos}

Trabalho Científico é fraco de Edição (Publica). Assina liga o trabalho aos seus autores, com pelo menos um autor (RF18). Cobre liga o trabalho à inscrição que o cobre. Cada trabalho tem no máximo uma inscrição de cobertura (RF19, RE9), e o trabalho que passa do limite da edição gera um item pela Sobretaxa. Termo é forte e cada instância é uma versão. Responde liga a inscrição à versão do termo apresentada, com o aceite e a data, então é possível saber qual versão valia em cada inscrição (RF23).

\begin{figure}[h!]
\centering
\merE
\caption{Parte E do modelo conceitual: trabalhos científicos e termos}
\label{fig:merE}
\end{figure}

\subsubsection*{Parte F: certificação e auditoria}

Certificado é forte e se liga à pessoa (Certifica), à edição (Emite) e, no certificado de atividade, à atividade (Atesta). Ele não depende da inscrição, porque organizadores e autores podem receber certificado sem estar inscritos (RF22, RE16). Auditoria é forte e Executa identifica a pessoa responsável pela operação, quando houver (RF25). Os losangos Certifica e Executa ficam à esquerda desta região, perto de Edição, e aparecem na Figura~\ref{fig:merfinal}. O recorte mostra as linhas que chegam deles a Certificado e Auditoria.

\begin{figure}[h!]
\centering
\merF
\caption{Parte F do modelo conceitual: certificados e auditoria}
\label{fig:merF}
\end{figure}

\FloatBarrier
\subsection{Dicionário das entidades}

A Tabela~\ref{tab:dicionario} define os atributos de cada entidade, com o tipo usado na implantação. Todas as 20 entidades têm um identificador \cod{cp_id_*}, o mesmo círculo cheio do diagrama. Nas entidades fracas e associativas, a identificação pela dona continua valendo: o identificador parcial, junto com a dona, é único. As chaves estrangeiras não aparecem, porque a ligação é o próprio relacionamento. Os domínios dos atributos de tipo e de situação estão na Seção~\ref{sec:restricoes}.

{\small
\begin{longtable}{@{}L{4.4cm} L{3.1cm} L{1.9cm} L{4.1cm}@{}}
\caption{Dicionário de atributos das entidades}\label{tab:dicionario}\\
\toprule
Atributo & Tipo & Obrigatório & Chave\\
\endfirsthead
\toprule
Atributo & Tipo & Obrigatório & Chave\\
\endhead
\bottomrule
\endfoot
\midrule
\multicolumn{4}{@{}l}{\textbf{1. Pessoa} (forte)}\\
\midrule
\cod{cp_id_pessoa} & \cod{bigint} & sim & identificador\\
\cod{nm_pessoa} & \cod{varchar(200)} & sim & \\
\cod{email_pessoa} & \cod{varchar(254)} & sim & único\\
\cod{pais_pessoa} & \cod{char(2)} & sim & \\
\cod{uf_pessoa} & \cod{char(2)} & não & \\
\cod{cpf_pessoa} & \cod{char(11)} & não & único\\
\cod{passaporte_pessoa} & \cod{varchar(30)} & não & \\
\cod{dt_nascimento} & \cod{date} & não & \\
\cod{ind_anonimizada} & \cod{boolean} & sim & \\
\midrule
\multicolumn{4}{@{}l}{\textbf{2. Evento} (forte)}\\
\midrule
\cod{cp_id_evento} & \cod{bigint} & sim & identificador\\
\cod{sg_evento} & \cod{varchar(20)} & sim & único\\
\cod{nm_evento} & \cod{varchar(200)} & sim & \\
\cod{tp_evento} & \cod{varchar(20)} & sim & \\
\cod{ds_evento} & \cod{varchar(4000)} & não & \\
\midrule
\multicolumn{4}{@{}l}{\textbf{3. Edição} (fraca, dona: Evento)}\\
\midrule
\cod{ano_edicao} & \cod{smallint} & sim & identificador parcial\\
\cod{nr_edicao} & \cod{int} & não & \\
\cod{dt_inicio_evento} & \cod{date} & sim & \\
\cod{dt_fim_evento} & \cod{date} & sim & \\
\cod{dt_inicio_inscricao} & \cod{timestamp} & sim & \\
\cod{dt_fim_inscricao} & \cod{timestamp} & sim & \\
\cod{modalidade_evento} & \cod{varchar(10)} & sim & \\
\cod{cidade_evento} & \cod{varchar(100)} & sim & \\
\cod{uf_evento} & \cod{char(2)} & não & \\
\cod{pais_evento} & \cod{char(2)} & sim & \\
\cod{nr_dias_reserva} & \cod{smallint} & sim & \\
\cod{nr_dias_limite_reembolso} & \cod{smallint} & sim & \\
\cod{pc_reembolso} & \cod{numeric(5,2)} & sim & \\
\cod{ind_exige_credenciamento} & \cod{boolean} & sim & \\
\cod{pc_frequencia_minima} & \cod{numeric(5,2)} & não & \\
\cod{nr_max_trabalhos_inscricao} & \cod{smallint} & não & \\
\cod{vl_taxa_trabalho_extra} & \cod{numeric(10,2)} & não & \\
\midrule
\multicolumn{4}{@{}l}{\textbf{4. Categoria de Inscrição} (fraca, dona: Edição)}\\
\midrule
\cod{nm_categoria} & \cod{varchar(100)} & sim & identificador parcial\\
\cod{ind_exige_comprovacao} & \cod{boolean} & sim & \\
\cod{ind_categoria_autor} & \cod{boolean} & sim & \\
\cod{ds_categoria} & \cod{varchar(4000)} & não & \\
\midrule
\multicolumn{4}{@{}l}{\textbf{5. Faixa de Inscrição} (fraca, dona: Edição)}\\
\midrule
\cod{nm_faixa} & \cod{varchar(60)} & sim & identificador parcial\\
\cod{dt_inicio_faixa} & \cod{timestamp} & sim & \\
\cod{dt_fim_faixa} & \cod{timestamp} & sim & \\
\midrule
\multicolumn{4}{@{}l}{\textbf{6. Preço} (associativa (Categoria × Faixa))}\\
\midrule
\cod{tp_associacao} & \cod{varchar(20)} & sim & identificador parcial\\
\cod{vl_inscricao} & \cod{numeric(10,2)} & sim & \\
\cod{vl_associacao} & \cod{numeric(10,2)} & não & \\
\midrule
\multicolumn{4}{@{}l}{\textbf{7. Atividade} (fraca, dona: Edição)}\\
\midrule
\cod{nm_atividade} & \cod{varchar(200)} & sim & identificador parcial\\
\cod{dt_inicio} & \cod{timestamp} & sim & identificador parcial\\
\cod{tp_atividade} & \cod{varchar(20)} & sim & \\
\cod{dt_fim} & \cod{timestamp} & sim & \\
\cod{ind_inscricao_adicional} & \cod{boolean} & sim & \\
\cod{nr_capacidade} & \cod{int} & não & \\
\cod{vl_atividade} & \cod{numeric(10,2)} & não & \\
\midrule
\multicolumn{4}{@{}l}{\textbf{8. Instituição} (forte)}\\
\midrule
\cod{cp_id_instituicao} & \cod{bigint} & sim & identificador\\
\cod{nm_instituicao} & \cod{varchar(200)} & sim & \\
\cod{sg_instituicao} & \cod{varchar(30)} & sim & \\
\cod{pais_instituicao} & \cod{char(2)} & sim & \\
\cod{ind_associada} & \cod{boolean} & sim & \\
\cod{cnpj_instituicao} & \cod{char(14)} & não & único\\
\cod{cidade_instituicao} & \cod{varchar(100)} & não & \\
\cod{uf_instituicao} & \cod{char(2)} & não & \\
\midrule
\multicolumn{4}{@{}l}{\textbf{9. Vínculo Institucional} (associativa (Pessoa × Instituição))}\\
\midrule
\cod{dt_inicio_vinculo} & \cod{date} & sim & identificador parcial\\
\cod{tp_vinculo} & \cod{varchar(20)} & sim & \\
\cod{dt_fim_vinculo} & \cod{date} & não & \\
\midrule
\multicolumn{4}{@{}l}{\textbf{10. Benefício} (fraca, dona: Instituição)}\\
\midrule
\cod{tp_beneficio} & \cod{varchar(20)} & sim & identificador parcial\\
\cod{dt_inicio_vigencia} & \cod{date} & sim & identificador parcial\\
\cod{dt_fim_vigencia} & \cod{date} & sim & \\
\cod{nr_limite_periodo} & \cod{int} & sim & \\
\cod{nr_limite_por_edicao} & \cod{int} & sim & \\
\cod{pc_desconto} & \cod{numeric(5,2)} & não & \\
\midrule
\multicolumn{4}{@{}l}{\textbf{11. Período de Associação} (fraca, dona: Pessoa)}\\
\midrule
\cod{dt_inicio_associacao} & \cod{date} & sim & identificador parcial\\
\cod{dt_fim_associacao} & \cod{date} & sim & \\
\midrule
\multicolumn{4}{@{}l}{\textbf{12. Inscrição} (forte)}\\
\midrule
\cod{cp_id_inscricao} & \cod{bigint} & sim & identificador\\
\cod{dt_inscricao} & \cod{timestamp} & sim & \\
\cod{status_inscricao} & \cod{varchar(25)} & sim & \\
\cod{dt_validade_preco} & \cod{timestamp} & sim & \\
\cod{dt_confirmacao} & \cod{timestamp} & não & \\
\cod{dt_cancelamento} & \cod{timestamp} & não & \\
\cod{dt_credenciamento} & \cod{timestamp} & não & \\
\midrule
\multicolumn{4}{@{}l}{\textbf{13. Inscrição em Atividade} (associativa (Inscrição × Atividade))}\\
\midrule
\cod{dt_inclusao} & \cod{timestamp} & sim & \\
\cod{status_inscricao_atividade} & \cod{varchar(20)} & sim & \\
\cod{dt_expira_reserva} & \cod{timestamp} & não & \\
\cod{dt_cancelamento} & \cod{timestamp} & não & \\
\cod{dt_presenca} & \cod{timestamp} & não & \\
\midrule
\multicolumn{4}{@{}l}{\textbf{14. Item Financeiro} (fraca, dona: Inscrição)}\\
\midrule
\cod{nr_item} & \cod{smallint} & sim & identificador parcial\\
\cod{tp_item} & \cod{varchar(20)} & sim & \\
\cod{vl_item} & \cod{numeric(10,2)} & sim & \\
\cod{status_item} & \cod{varchar(15)} & sim & \\
\cod{dt_criacao} & \cod{timestamp} & sim & \\
\midrule
\multicolumn{4}{@{}l}{\textbf{15. Pagamento} (fraca, dona: Inscrição)}\\
\midrule
\cod{dt_criacao} & \cod{timestamp} & sim & identificador parcial\\
\cod{tp_pagamento} & \cod{varchar(15)} & sim & \\
\cod{vl_pagamento} & \cod{numeric(10,2)} & sim & \\
\cod{status_pagamento} & \cod{varchar(15)} & sim & \\
\cod{dt_confirmacao} & \cod{timestamp} & não & \\
\cod{dt_vencimento} & \cod{date} & não & \\
\cod{cd_transacao_externa} & \cod{varchar(100)} & não & único\\
\cod{vl_reembolso} & \cod{numeric(10,2)} & não & \\
\cod{dt_reembolso} & \cod{timestamp} & não & \\
\midrule
\multicolumn{4}{@{}l}{\textbf{16. Documento de Comprovação} (fraca, dona: Inscrição)}\\
\midrule
\cod{dt_envio} & \cod{timestamp} & sim & identificador parcial\\
\cod{tp_documento} & \cod{varchar(30)} & sim & \\
\cod{status_documento} & \cod{varchar(10)} & sim & \\
\cod{ds_caminho_arquivo} & \cod{varchar(500)} & sim & \\
\cod{dt_analise} & \cod{timestamp} & não & \\
\cod{ds_motivo_rejeicao} & \cod{varchar(500)} & não & \\
\midrule
\multicolumn{4}{@{}l}{\textbf{17. Trabalho Científico} (fraca, dona: Edição)}\\
\midrule
\cod{cd_trabalho_externo} & \cod{varchar(50)} & sim & identificador parcial\\
\cod{titulo_trabalho} & \cod{varchar(500)} & sim & \\
\cod{tp_trabalho} & \cod{varchar(20)} & sim & \\
\cod{status_trabalho} & \cod{varchar(15)} & sim & \\
\midrule
\multicolumn{4}{@{}l}{\textbf{18. Termo} (forte)}\\
\midrule
\cod{cp_id_termo} & \cod{bigint} & sim & identificador\\
\cod{nm_termo} & \cod{varchar(100)} & sim & \\
\cod{nr_versao} & \cod{varchar(10)} & sim & \\
\cod{ds_termo} & \cod{varchar(4000)} & sim & \\
\cod{dt_inicio_vigencia} & \cod{timestamp} & sim & \\
\cod{dt_fim_vigencia} & \cod{timestamp} & não & \\
\midrule
\multicolumn{4}{@{}l}{\textbf{19. Certificado} (forte)}\\
\midrule
\cod{cp_id_certificado} & \cod{bigint} & sim & identificador\\
\cod{cd_validacao} & \cod{varchar(64)} & sim & único\\
\cod{tp_certificado} & \cod{varchar(20)} & sim & \\
\cod{dt_emissao} & \cod{timestamp} & sim & \\
\cod{status_certificado} & \cod{varchar(10)} & sim & \\
\midrule
\multicolumn{4}{@{}l}{\textbf{20. Auditoria} (forte)}\\
\midrule
\cod{cp_id_auditoria} & \cod{bigint} & sim & identificador\\
\cod{nm_tabela} & \cod{varchar(60)} & sim & \\
\cod{cd_registro} & \cod{bigint} & sim & \\
\cod{tp_operacao} & \cod{varchar(10)} & sim & \\
\cod{dt_operacao} & \cod{timestamp} & sim & \\
\cod{ds_valor_anterior} & \cod{varchar(4000)} & não & \\
\cod{ds_valor_novo} & \cod{varchar(4000)} & não & \\
\end{longtable}

}

\subsection{Relacionamentos}

A Tabela~\ref{tab:relacionamentos} lista os 30 relacionamentos. A cardinalidade está no formato (junto a A) e (junto a B).

{\small
\begin{longtable}{@{}l L{3.0cm} L{2.3cm} L{2.3cm} L{1.8cm} L{2.6cm}@{}}
\caption{Relacionamentos do modelo conceitual}\label{tab:relacionamentos}\\
\toprule
\# & Relacionamento & A & B & Card. & Atributos\\
\midrule
\endfirsthead
\toprule
\# & Relacionamento & A & B & Card. & Atributos\\
\midrule
\endhead
\bottomrule
\endfoot
R01 & Realiza (identificador) & Evento & Edição & (1,1) (1,n) & \\
R02 & Agrega (autorrelac.) & Edição agregadora & Edição agregada & (0,1) (0,n) & \\
R03 & Oferece (identificador) & Edição & Categoria & (1,1) (1,n) & \\
R04 & Estabelece (identificador) & Edição & Faixa & (1,1) (1,n) & \\
R05 & Programa (identificador) & Edição & Atividade & (1,1) (0,n) & \\
R06 & Publica (identificador) & Edição & Trabalho & (1,1) (0,n) & \\
R07 & Preço (associativa) & Categoria & Faixa & (1,n) (1,n) & tipo de associação, valores\\
R08 & Enquadra & Preço & Inscrição & (1,1) (0,n) & \\
R09 & Efetua & Pessoa & Inscrição & (1,1) (0,n) & \\
R10 & Facilita & Benefício & Inscrição & (0,1) (0,n) & \cod{dt_uso_beneficio}\\
R11 & Vínculo Institucional (associativa) & Pessoa & Instituição & (0,n) (0,n) & tipo e período\\
R12 & Recebe (identificador) & Instituição & Benefício & (1,1) (0,n) & \\
R13 & Associa-se (identificador) & Pessoa & Período & (1,1) (0,n) & \\
R14 & Origina & Inscrição & Período & (0,1) (0,1) & \\
R15 & Inscrição em Atividade (associativa) & Inscrição & Atividade & (0,n) (0,n) & inclusão, situação, reserva, cancelamento, presença\\
R16 & Compõe (identificador) & Inscrição & Item Financeiro & (1,1) (1,n) & \\
R17 & Tarifa & Inscrição em Atividade & Item Financeiro & (0,1) (0,n) & \\
R18 & Sobretaxa & Trabalho & Item Financeiro & (0,1) (0,n) & \\
R19 & Transaciona (identificador) & Inscrição & Pagamento & (1,1) (0,n) & \\
R20 & Garante (identificador) & Inscrição & Documento & (1,1) (0,n) & \\
R21 & Analisa & Pessoa & Documento & (0,1) (0,n) & \\
R22 & Cobre & Inscrição & Trabalho & (0,1) (0,n) & \\
R23 & Assina & Pessoa & Trabalho & (1,n) (0,n) & \\
R24 & Responde & Inscrição & Termo & (0,n) (0,n) & \cod{ind_aceite}, \cod{dt_resposta}\\
R25 & Organiza & Pessoa & Edição & (1,n) (0,n) & \cod{tp_papel}\\
R26 & Restringe & Atividade & Categoria & (0,n) (0,n) & \\
R27 & Certifica & Pessoa & Certificado & (1,1) (0,n) & \\
R28 & Emite & Edição & Certificado & (1,1) (0,n) & \\
R29 & Atesta & Atividade & Certificado & (0,1) (0,n) & \\
R30 & Executa & Pessoa & Auditoria & (0,1) (0,n) & \\
\end{longtable}
}

\subsection{Modelo lógico}
\label{sec:logico}

O modelo lógico foi gerado no brModelo pela conversão do modelo conceitual e depois revisado à mão, tabela por tabela, contra a DDL do banco. Ele tem 24 tabelas: uma para cada uma das 20 entidades e uma para cada relacionamento N:N sem entidade própria (Assina, Responde, Organiza e Restringe). A Seção~\ref{sec:traducao} detalha as regras de tradução.

A conversão automática teve quatro problemas, que o grupo corrigiu. O relacionamento Recebe, que é 1:N, tinha virado tabela e passou a ser a chave estrangeira \cod{ce_instituicao} em Benefício. Responde, que é N:N com atributos, não tinha gerado tabela e foi criada com \cod{ind_aceite} e \cod{dt_resposta}. Vínculo Institucional, Preço e Inscrição em Atividade usavam as chaves estrangeiras como chave primária, o que impediria, por exemplo, a mesma pessoa voltar à mesma instituição em outro período. As três passaram a ter chave própria (\cod{cp_id_*}), como no banco. Por fim, as chaves estrangeiras que faltavam (como as de Pagamento, Item Financeiro e Documento para Inscrição) foram incluídas, e todas foram renomeadas para o padrão \cod{ce_*} da Seção~\ref{sec:nomes}.

Para ler as figuras: cada retângulo é uma tabela, PK marca a chave primária e FK marca a chave estrangeira. Cada linha liga uma chave estrangeira à tabela que ela referencia, com a cardinalidade nas pontas. As 37 chaves estrangeiras do banco aparecem no modelo. Só \cod{ce_edicao_agregadora}, em Edição, não tem linha desenhada, porque aponta para a própria tabela e o brModelo não desenha ligação de uma tabela com ela mesma.

O brModelo oferece um conjunto pequeno de tipos, então três diferenças com o banco ficam no diagrama. Valores monetários e percentuais aparecem como \cod{FLOAT}, e no banco são \cod{numeric}. Chaves aparecem como \cod{INT}, e no banco são \cod{bigint}. Data com hora aparece como \cod{DATETIME}, que no PostgreSQL é \cod{timestamp}. O brModelo também força o tipo \cod{INT} em toda coluna marcada como PK. Por isso \cod{tp_papel}, que faz parte da chave de Organiza e no banco é \cod{varchar(20)}, aparece como \cod{INT}, com uma nota ao lado indicando o tipo real. Os tipos exatos estão no dicionário (Tabela~\ref{tab:dicionario}) e na DDL.

A Figura~\ref{fig:loggeral} mostra o modelo completo, e as Figuras~\ref{fig:logA} a~\ref{fig:logG} ampliam cada região, também em formato vetorial.

\begin{landscape}
\begin{figure}[p]
\centering
\loggeral
\caption{Modelo lógico do EventOS}
\label{fig:loggeral}
\end{figure}
\end{landscape}

\begin{figure}[h!]
\centering
\logA
\caption{Parte A do modelo lógico: Pessoa, Período de Associação, Assina, Auditoria e Organiza}
\label{fig:logA}
\end{figure}

\begin{figure}[h!]
\centering
\logB
\caption{Parte B do modelo lógico: Vínculo Institucional, Instituição, Benefício e Documento de Comprovação}
\label{fig:logB}
\end{figure}

\begin{figure}[h!]
\centering
\logC
\caption{Parte C do modelo lógico: Evento e Edição}
\label{fig:logC}
\end{figure}

\begin{figure}[h!]
\centering
\logD
\caption{Parte D do modelo lógico: Faixa de Inscrição, Preço, Categoria de Inscrição e Restringe}
\label{fig:logD}
\end{figure}

\begin{figure}[h!]
\centering
\logE
\caption{Parte E do modelo lógico: Certificado e Atividade}
\label{fig:logE}
\end{figure}

\begin{figure}[h!]
\centering
\logF
\caption{Parte F do modelo lógico: Inscrição, Trabalho Científico, Pagamento, Termo e Responde}
\label{fig:logF}
\end{figure}

\begin{figure}[h!]
\centering
\logG
\caption{Parte G do modelo lógico: Inscrição em Atividade e Item Financeiro}
\label{fig:logG}
\end{figure}

\FloatBarrier
